\documentclass[a4paper,12pt,twoside]{ctexart}
\usepackage[dvipsnames]{xcolor}
\usepackage{geometry}
\usepackage{tcolorbox}
\usepackage{tikz}
\usepackage{amssymb}
\usepackage{graphicx}

\definecolor{inkblack}{rgb}{0.08,0.08,0.08}
\geometry{a4paper, left=3.18cm, right=3.18cm, top=2.54cm, bottom=2.54cm}

\usepackage{fancyhdr}
\pagestyle{fancy}
\fancyhf{}
\renewcommand{\headrulewidth}{0pt}
\fancyfoot[RO]{\makebox[0pt][l]{\hspace*{28.2pt}\raisebox{2.81pt}{\rmfamily\fontsize{10.5bp}{10.5bp}\selectfont\thepage}}}
\fancyfoot[LE]{\makebox[0pt][r]{\hspace*{33.45pt}\raisebox{2.81pt}{\rmfamily\fontsize{10.5bp}{10.5bp}\selectfont\thepage}}}

\newCJKfontfamily\fzdbs[Path=./]{FZDBS.ttf}
\newCJKfontfamily\fzssk[Path=./]{FZSSK.ttf}
\AtBeginDocument{\color{inkblack}}

\newcommand{\lawtitle}[2]{
    \vspace*{0.507cm}
    \begin{center}
        {\fzdbs\fontsize{22bp}{22bp}\selectfont #1}\par
        \vspace{10.37bp}
        {\fzdbs\fontsize{18bp}{18bp}\selectfont #2}\par
    \end{center}
    \vspace{18.4bp}
}

\newtcolorbox{notesection}[1][]{
    width=467.72bp, sharp corners, colback=white, colframe=inkblack, boxrule=0.5pt,
    fontupper=\color{inkblack}\fontsize{10.5bp}{10.5bp}\selectfont,
    enlarge left by=-19.28bp, before=\vspace*{-13.39bp}\noindent, after=\par,
    boxsep=0pt, left=4bp, top=4.69bp, bottom=4.69bp, right=0pt, #1
}

\newcommand{\lawpar}{\par\vspace*{-4.86bp}}

\newcommand{\subtitlesection}[1]{%
    \begin{tcolorbox}[
        width=467.72bp, sharp corners, colback=white, colframe=inkblack, boxrule=0.5pt,
        boxsep=0pt, left=0pt, right=0pt, top=4bp, bottom=4bp, halign=center,
        fontupper=\fzdbs\fontsize{15bp}{15bp}\selectfont\color{inkblack},
        enlarge left by=-19.28bp,
        before={\nointerlineskip\vspace{-0.5pt}\noindent}, after=\vspace{0pt}
    ]
    #1
    \end{tcolorbox}%
}

\newcommand{\subtitlenotesection}[2]{%
    \begin{tcolorbox}[
        width=467.72bp, sharp corners, colback=white, colframe=inkblack, boxrule=0.5pt,
        boxsep=0pt, left=0pt, right=0pt, top=4bp, bottom=4bp, halign=center,
        enlarge left by=-19.28bp,
        before={\nointerlineskip\vspace{-0.5pt}\noindent}, after=\vspace{0pt}
    ]
    {\fzdbs\fontsize{15bp}{15bp}\selectfont\color{inkblack} #1}\par
    {\songti\fontsize{10.5bp}{14bp}\selectfont\color{inkblack}\bfseries #2}
    \end{tcolorbox}%
}

\newcommand{\pairsectionbase}[3]{%
    \begin{tcolorbox}[
        width=467.72bp, sharp corners, colback=white, colframe=inkblack, boxrule=0.5pt,
        enlarge left by=-19.28bp,
        before={\nointerlineskip\vspace{-0.5pt}\noindent}, after=\par,
        boxsep=0pt, left=0pt, right=0pt, top=0pt, bottom=0pt
    ]
    \begin{minipage}[t]{112.9bp}%
        \vspace{4bp}%
        \songti\fontsize{10.5bp}{12bp}\selectfont
        #1#2%
        \vspace{4bp}%
    \end{minipage}%
    \hspace{0pt}%
    \vrule width 0.5pt%
    \hspace{4bp}%
    \begin{minipage}[c]{346.32bp}%
        \vspace{4bp}%
        \songti\mdseries\fontsize{10.5bp}{17bp}\selectfont%
        \baselineskip=17bp%
        \setlength{\parindent}{0pt}%
        \color{inkblack}%
        #3%
        \vspace{4bp}%
    \end{minipage}%
    \end{tcolorbox}%
}

\newcommand{\pairsection}[2]{%
    \pairsectionbase{\centering}{#1}{#2}%
}

\newcommand{\pairsectionleft}[2]{%
    \pairsectionbase{\raggedright\fontsize{10.5bp}{10.5bp}\selectfont}{#1}{#2}%
}

\begin{document}

% ===== 第 1 页 =====
\lawtitle{民事起诉状}{（房屋租赁合同纠纷）}

\begin{notesection}
\hfuzz=50pt
\color{inkblack}\heiti\fontsize{10.5bp}{16bp}\selectfont 说明：\par
\vspace*{-0.2bp}
\songti\mdseries\fontsize{10.5bp}{16bp}\selectfont
\setlength{\parindent}{20.7bp}
\sloppy
\setlength{\parskip}{0pt}
\vspace*{-5.1bp}为了方便您更好地参加诉讼，保护您的合法权利，请填写本表。
\lawpar
1.起诉时需向人民法院提交证明您身份的材料，如身份证复印件、营业执照复印件等。
\lawpar
2.本表所列内容是您提起诉讼以及人民法院查明案件事实所需，请务必如实填写。
\lawpar
3.本表有些内容可能与您的案件无关，您认为与案件无关的项目可以填"无"或不填；对于本\\[-4.86bp]
表中勾选项可以在对应项打"√"；您认为另有重要内容需要列明的，可以另附页填写。
\lawpar
4.本表 word 电子版填写时，相关栏目可复制粘贴或扩容，但不得改变要素内容、格式设置。例\\[-4.86bp]
如，多原告、多被告或多委托诉讼代理人等情况，可根据实际情况复制粘贴；需填写文字较多时，\\[-4.86bp]
可根据实际对栏目进行扩容等。
\lawpar
$\bigstar$特别提示$\bigstar$\\[-4.86bp]
\hspace*{20.7bp}诉讼参加人应遵守诚信原则如实认真填写表格。\\[-4.86bp]
\hspace*{20.7bp}如果诉讼参加人违反有关规定，虚假诉讼、恶意诉讼、滥用诉权，人民法院将视违法情形依法\\[-4.86bp]
追究责任。
\end{notesection}

\subtitlesection{当事人信息}

\pairsection{原告\\（自然人）}{
姓名：\\
性别：男$\square$\hspace{1em}女$\square$\\
出生日期：\hspace{2em}年\hspace{2em}月\hspace{2em}日\hspace{4em}民族：\\
工作单位：\hspace{6em}职务：\hspace{4em}联系电话：\\
住所地（户籍所在地）：\\
经常居住地：\\
证件类型：\\
证件号码：
}

\pairsection{原告\\（法人、非法人组织）}{
名称：\\
住所地（主要办事机构所在地）：\\
注册地 / 登记地：\\
法定代表人 / 负责人：\hspace{4em}职务：\hspace{4em}联系电话：\\
统一社会信用代码：\\
类型：有限责任公司$\square$\hspace{1em}股份有限公司$\square$\hspace{1em}上市公司$\square$\\
\hspace*{36bp}其他企业法人$\square$\hspace{1em}事业单位$\square$\hspace{1em}社会团体$\square$\hspace{1em}基金会$\square$\\
\hspace*{36bp}社会服务机构$\square$\hspace{1em}机关法人$\square$\hspace{1em}农村集体经济组织法人$\square$\\
\hspace*{36bp}城镇农村的合作经济组织法人$\square$\hspace{1em}基层群众性自治组织法人$\square$\\
\hspace*{36bp}个人独资企业$\square$\hspace{1em}合伙企业$\square$\hspace{1em}不具有法人资格的专业服务机构$\square$\\
所有制性质：国有$\square$（控股$\square$\hspace{1em}参股$\square$）\hspace{1em}民营$\square$\hspace{1em}其他\rule{60bp}{0.5pt}
}

\newpage
% ===== 第 2 页 =====

\pairsection{委托诉讼代理人}{
有$\square$\\
\hspace*{21bp}姓名：\\
\hspace*{21bp}\makebox[0.33\linewidth][l]{单位：}\makebox[0.33\linewidth][l]{职务：}\makebox[0.33\linewidth][l]{联系电话：}\\
\hspace*{21bp}代理权限：一般授权$\square$\hspace{1em}特别授权$\square$\hspace{1em}\rule{80bp}{0.5pt}\\
无$\square$
}

\pairsection{被告\\（自然人）}{
姓名：\\
性别：男$\square$\hspace{1em}女$\square$\\
出生日期：\hspace{2em}年\hspace{2em}月\hspace{2em}日\hspace{4em}民族：\\
工作单位：\hspace{6em}职务：\hspace{4em}联系电话：\\
住所地（户籍所在地）：\\
经常居住地：\\
证件类型：\\
证件号码：
}

\pairsection{被告\\（法人、非法人组织）}{
名称：\\
住所地（主要办事机构所在地）：\\
注册地 / 登记地：\\
法定代表人 / 负责人：\hspace{4em}职务：\hspace{4em}联系电话：\\
统一社会信用代码：\\
类型：有限责任公司$\square$\hspace{1em}股份有限公司$\square$\hspace{1em}上市公司$\square$\\
\hspace*{36bp}其他企业法人$\square$\hspace{1em}事业单位$\square$\hspace{1em}社会团体$\square$\hspace{1em}基金会$\square$\\
\hspace*{36bp}社会服务机构$\square$\hspace{1em}机关法人$\square$\hspace{1em}农村集体经济组织法人$\square$\\
\hspace*{36bp}城镇农村的合作经济组织法人$\square$\hspace{1em}基层群众性自治组织法人$\square$\\
\hspace*{36bp}个人独资企业$\square$\hspace{1em}合伙企业$\square$\hspace{1em}不具有法人资格的专业服务机构$\square$\\
所有制性质：国有$\square$（控股$\square$\hspace{1em}参股$\square$）\hspace{1em}民营$\square$\hspace{1em}其他\rule{60bp}{0.5pt}
}

\pairsection{第三人\\（自然人）}{
姓名：\\
性别：男$\square$\hspace{1em}女$\square$\\
出生日期：\hspace{2em}年\hspace{2em}月\hspace{2em}日\hspace{4em}民族：\\
工作单位：\hspace{6em}职务：\hspace{4em}联系电话：\\
住所地（户籍所在地）：\\
经常居住地：\\
证件类型：\\
证件号码：
}

\newpage
% ===== 第 3 页 =====

\pairsection{第三人\\（法人、非法人组织）}{
名称：\\
住所地（主要办事机构所在地）：\\
注册地 / 登记地：\\
法定代表人 / 负责人：\hspace{4em}职务：\hspace{4em}联系电话：\\
统一社会信用代码：\\
类型：有限责任公司$\square$\hspace{1em}股份有限公司$\square$\hspace{1em}上市公司$\square$\\
\hspace*{36bp}其他企业法人$\square$\hspace{1em}事业单位$\square$\hspace{1em}社会团体$\square$\hspace{1em}基金会$\square$\\
\hspace*{36bp}社会服务机构$\square$\hspace{1em}机关法人$\square$\hspace{1em}农村集体经济组织法人$\square$\\
\hspace*{36bp}城镇农村的合作经济组织法人$\square$\hspace{1em}基层群众性自治组织法人$\square$\\
\hspace*{36bp}个人独资企业$\square$\hspace{1em}合伙企业$\square$\hspace{1em}不具有法人资格的专业服务机构$\square$\\
所有制性质：国有$\square$（控股$\square$\hspace{1em}参股$\square$）\hspace{1em}民营$\square$\hspace{1em}其他\rule{60bp}{0.5pt}
}

\subtitlenotesection{诉讼请求}{（原告为出租方时，填写第 1 项、第 2 项、第 5 项、第 6 项、第 7 项；原告为承租方时，\\填写第 3 项、第 8 项；第 4 项、第 9 项至第 12 项为共同项）}

\begin{tcolorbox}[
    width=467.72bp, sharp corners, colback=white, colframe=inkblack, boxrule=0.5pt,
    enlarge left by=-19.28bp,
    before={\nointerlineskip\vspace{-0.5pt}\noindent}, after=\par,
    boxsep=0pt, left=4bp, right=0pt, top=4bp, bottom=4bp
]
\songti\mdseries\fontsize{10.5bp}{17bp}\selectfont%
\baselineskip=17bp\parskip=0pt\lineskiplimit=-\maxdimen%
\setlength{\parindent}{0pt}\color{inkblack}%
（可完整表述诉讼请求；为方便、准确梳理要点，相关内容请在下方要素式表格中填写）\\[51bp]
\mbox{}
\end{tcolorbox}

\pairsectionleft{1. 支付租金（元）}{
到期未付租金\hspace{4em}元（人民币，下同；如外币需特别注明）\\
明细：
}

\pairsectionleft{2. 迟延支付租金的利息\\（违约金）}{
截至\hspace{2em}年\hspace{2em}月\hspace{2em}日止，迟延支付租金的利息\hspace{4em}元、违约\\
金\hspace{4em}元，自\hspace{4em}之后的逾期利息、违约金，以\hspace{4em}元为基数按\\
照\hspace{4em}标准计算；\\
计算方式：\\
是否请求支付至实际清偿之日止：是$\square$\hspace{1em}否$\square$\\
明细：
}

\pairsectionleft{3. 交付房屋}{
是$\square$\hspace{1em}否$\square$
}

\pairsectionleft{4. 请求解除合同}{
是$\square$\hspace{1em}确认合同于\hspace{2em}年\hspace{2em}月\hspace{2em}日解除\\
否$\square$
}

\pairsectionleft{5. 返还租赁物，并赔偿\\因解除合同而受到的\\损失}{
是$\square$\hspace{1em}内容：\\
否$\square$
}

\pairsectionleft{6. 支付房屋占有使用费}{
是$\square$\hspace{1em}内容：\\
否$\square$
}

\pairsectionleft{7. 支付水电费等费用}{
是$\square$\hspace{1em}内容：\\
否$\square$
}

\pairsectionleft{8. 返还押金}{
是$\square$\hspace{1em}金额：\\
否$\square$
}

\pairsectionleft{9. 是否主张实现债权的\\费用}{
是$\square$\hspace{1em}内容：\\
否$\square$
}

\pairsectionleft{10. 是否主张诉讼费用}{
是$\square$\\
否$\square$
}

\pairsectionleft{11. 其他请求}{
\mbox{}
}

\pairsectionleft{12. 标的总额}{
\mbox{}
}

\subtitlesection{约定管辖和诉前保全}

\pairsectionleft{1. 有无仲裁、法院管辖\\约定}{
有$\square$\hspace{4em}合同条款及内容：\\
无$\square$
}

\pairsectionleft{2. 是否已经诉前保全}{
是$\square$\hspace{1em}保全法院：\hspace{4em}保全时间：\\
\hspace*{21bp}保全案号：\\
否$\square$\\
（如申请诉讼保全，请另行提交诉讼保全申请及相关材料）
}

\newpage
% ===== 第 4 页 =====

\subtitlesection{事实与理由}

\begin{tcolorbox}[
    width=467.72bp, sharp corners, colback=white, colframe=inkblack, boxrule=0.5pt,
    enlarge left by=-19.28bp,
    before={\nointerlineskip\vspace{-0.5pt}\noindent}, after=\par,
    boxsep=0pt, left=4bp, right=0pt, top=4bp, bottom=4bp
]
\songti\mdseries\fontsize{10.5bp}{17bp}\selectfont%
\baselineskip=17bp\parskip=0pt\lineskiplimit=-\maxdimen%
\setlength{\parindent}{0pt}\color{inkblack}%
（可完整表述纠纷涉及的事实与理由；为方便、准确梳理要点，相关内容请在下方要素式表格中填写）\\[51bp]
\mbox{}
\end{tcolorbox}

\pairsectionleft{1. 合同的签订情况（名\\称、编号、签订时间、\\地点等）}{
\mbox{}
}

\pairsectionleft{2. 签订主体}{
出租人：\\
承租人：
}

\pairsectionleft{3. 租赁标的物情况（坐\\落位置、面积、产权\\情况等）}{
\mbox{}
}

\pairsectionleft{4. 合同约定的租赁期限}{
自\hspace{2em}年\hspace{2em}月\hspace{2em}日起至\hspace{2em}年\hspace{2em}月\hspace{2em}日止
}

\pairsectionleft{5. 合同约定的租金及支\\付方式}{
租金：\hspace{4em}元/月；总价\hspace{4em}元：\\
以现金$\square$\hspace{1em}转账$\square$\hspace{1em}票据$\square$（写明票据类型）\hspace{1em}其他$\square$\hspace{1em}方式\\
一次性$\square$\hspace{1em}分期$\square$\hspace{1em}支付\\
分期方式：
}

\pairsectionleft{6. 其他费用约定（物业\\费、水电燃气费用等）}{
出租人负担：\\
承租人负担：
}

\pairsectionleft{7. 合同约定的违约责任}{
\mbox{}
}

\pairsectionleft{8. 是否约定合同解除的\\条件}{
是$\square$\hspace{1em}具体内容：\\
否$\square$
}

\pairsectionleft{9. 租赁物交付时间}{
于\hspace{2em}年\hspace{2em}月\hspace{2em}日交付租赁物
}

\newpage
% ===== 第 5 页 =====

\pairsectionleft{10. 押金约定情况}{
有$\square$\hspace{2em}押金数额：\hspace{4em}，\hspace{2em}年\hspace{2em}月\hspace{2em}日已支付押金。\\
无$\square$
}

\pairsectionleft{11. 租金支付情况}{
自\hspace{2em}年\hspace{2em}月\hspace{2em}日至\hspace{2em}年\hspace{2em}月\hspace{2em}日，按约定交纳租金，已付租\\
金\hspace{4em}元，逾期但已付租金\hspace{4em}元\\
明细：
}

\pairsectionleft{12. 逾期未付租金情况}{
自\hspace{2em}年\hspace{2em}月\hspace{2em}日起开始欠付租金，截至\hspace{2em}年\hspace{2em}月\hspace{2em}日，欠付\\
租金\hspace{4em}元
}

\pairsectionleft{13. 其他需要说明的内\\容（可另附页）}{
\mbox{}
}

\pairsectionleft{14. 请求依据}{
合同约定：\\
法律规定：
}

\pairsectionleft{15. 证据清单（可另附\\页）}{
\mbox{}
}

\subtitlesection{对纠纷解决方式的意愿}

\pairsection{是否了解调解作为非诉\\讼纠纷解决方式，能\\及时、高效、低成本、\\不伤和气地解决纠纷}{
了解$\square$\hspace{1em}不了解$\square$
}

\pairsection{是否了解先行调解解\\决纠纷的好处}{
1. 立案后选择先行调解的，可以很快启动调解程序。如不同意调解，法院将\\
依程序开庭审理案件，但可能需要经过较长一段时间的排期等待，且审理、\\
执行周期相对较长。\\
了解$\square$\hspace{1em}不了解$\square$\\
2. 选择先行调解，调解成功且自动履行的免交诉讼费用，申请司法确认的不\\
交纳诉讼费用，要求出具调解书的减半交纳诉讼费用。\\
了解$\square$\hspace{1em}不了解$\square$\\
3. 首次调解不成功，但仍有继续调解意愿的，可以选择更换调解组织和调解\\
员再进行调解。调解无法达成一致意见的，法院将依程序排期开庭。\\
了解$\square$\hspace{1em}不了解$\square$\\
4. 依照法律规定，调解具有保密性要求，调解过程不公开，调解协议未经当\\
事人同意不得公开。\\
了解$\square$\hspace{1em}不了解$\square$\\
5. 调解达成的协议具有法律效力，可以依照法律规定申请司法确认，具有强\\
制执行效力。\\
了解$\square$\hspace{1em}不了解$\square$
}

\newpage
% ===== 第 6 页 =====

\pairsection{是否考虑先行调解}{
是$\square$\\
否$\square$\\
暂不确定，想要了解更多内容$\square$
}

\vspace{10bp}
\noindent\hspace*{248bp}{\fzdbs\fontsize{15bp}{22bp}\selectfont 具状人（签字、盖章）：}\par
\noindent\hspace*{248bp}{\fzdbs\fontsize{15bp}{22bp}\selectfont 日期：}

\end{document}
